\subsubsection{Cadenas, recuperación y límites de integración (CEN-20)}
\label{sec:sd4-cadenas-cen20}

La selección conserva cuarenta CSB \texttt{HRL12540W RA260131} por cadena y dos cadenas por UPS. El máximo de carga de la ficha es 54 A \emph{por bloque}, no por conjunto de ochenta. A 60 A agregados, el reparto ideal es 30 A por cadena; ni la igualdad de cables ni la conexión en paralelo garantizan ese reparto en todo estado. Si una cadena queda aislada, 60 A en la restante exceden en 6 A el límite CSB. La ficha identifica 25\,$^{\circ}$C como temperatura nominal y carga/almacenamiento de $-15$ a 40\,$^{\circ}$C; ese intervalo no demuestra disponibilidad de 54 A a todos sus extremos ni publica una curva de reducción térmica. Los 30 A son reparto nominal, no capacidad garantizada tras envejecimiento o calentamiento. El perfil autorizado y la compensación térmica se verifican en la banda del proyecto. La supervisión mide corriente de cada cadena, tensión de bloques y temperatura; una cadena retirada no conserva el permiso de recarga del conjunto. Se exige límite autorizado para el estado de una cadena y, mientras no esté individualizado, se inhibe la recarga de esa rama y se declara reserva degradada. La corriente de cortocircuito o los 200 A de la protección DC no son un ajuste de protección de carga a 54 A \parencites[p.~1]{csbHRL12540WRA260131}[pp.~49--51 y 98]{p7-eaton-93t-guia}.

\textbf{Interfaz que puede documentarse.} CN12 recibe alarmas externas mediante contacto aislado y cable trenzado; la guía remite su configuración a Eaton. Los tres juegos de entradas configurables y el mando local de cargador no publican una función específica de inhibición selectiva utilizable como interfaz ya autorizada. Se cablean por separado alarma de caudal, temperatura y estado de cadenas; no se inyecta tensión de potencia en CN12. La apertura externa de AC del rectificador se estudia únicamente como retirada de la fuente de carga: con el modo Frequency Converter y bypass bloqueado, la UPS pasa a batería y se consume su reserva. No equivale a mantener el inversor indefinidamente desde red sin cargar, ni a restaurar una autonomía agotada. La secuencia requiere comprobar el camino sano, transferir las cargas que lo admitan, retirar recarga y aislar el banco afectado según seguridad. Un banco aislado puede seguir liberando gas; extracción y comprobación local permanecen habilitadas. El circuito externo, variante de maniobra, fuente de control, retorno y aceptación Eaton deben seleccionarse juntos; no se atribuye ese circuito a una función documentada de CN12 o REPO \parencites[secc.~5.4.2, 5.4.7, 5.5.2 y 6.2]{p7-eaton-93t-guia}{csbManualJunio2025}.

\textbf{Recuperación y cortes sucesivos.} Con 60 kW de salida durante una hora y eficiencia propia de descarga 0,94, la energía DC retirada sería al menos $60/0,94=63,8298$ kWh en ese modelo. A 588 V y 60 A, la potencia máxima DC es 35,28 kW; aun ignorando absorción final, pérdidas de batería y restricciones de entrada, restituir esa energía exige al menos 1,8092 h. Para treinta minutos, el límite inferior ideal es 0,9046 h. Son cotas inferiores para detectar plazos imposibles, no tiempos ofrecidos de recarga ni capacidades nominales en Ah. A plena salida, CEN-4 limita por cribado la corriente a 16,3806 A si se supone FP 0,99 y rendimiento 0,94: el máximo de 60 A no permite ofrecer recuperación a plena carga. La autorización de disponibilidad utiliza descarga--recarga--nuevo corte de cada rama, temperaturas y estado de todas sus cadenas; no basta una tensión de flotación ni el aviso de red restablecida.

\textbf{Entradas únicas.} La hoja de continuidad individual conserva NVR 200 W/L3, dos terminaciones WAN de 75 W/L2, radio 300 W/L2, dos motores de 600 W/L3, dieciséis juegos móviles de 50 W y conversión DC/PoE. Estas partidas no se retiran del parque ni se suman nuevamente por seleccionar interfaces. Cada una debe disponer de tensión y frecuencia nominales, peor interrupción y tensión mínima admisible, sostenimiento documentado de su fuente, transferencia seleccionada y protección. El ensayo de una transferencia comprueba pérdida de entrada y de fuente interna separadamente: una PSU única averiada no se resuelve con cambiar su entrada AC. Para carga DC se exige conversión A/B desacoplada con comportamiento ante cortocircuito; para AC, el peor intervalo de la transferencia debe ser menor que el sostenimiento documentado del receptor. Sin esos datos no se asigna un STS genérico ni se declara continuidad acreditada por duplicar cables. Los equipos de control de campo mantienen estado seguro y liberación de emergencia durante las maniobras.

\subsubsection{Extracción seleccionada, fases y estados (ECI-20)}
\label{sec:sd4-extraccion-eci20}

ENV-19 selecciona ocho Systemair \texttt{EX 180A-4}, referencia 135279, distribuidos en cuatro trenes de dos motores: un tren activo y otro de reserva por banco. La partida por motor es 0,248 kW y 0,73 A, a 400 V trifásicos/50 Hz. Cuatro activos suman 0,992 kW y $4\sqrt{3}(400)(0,73)/1000=2,0230$ kVA; se conserva la corriente publicada y no se divide la potencia por un FP inventado. Ocho motores simultáneos suman 1,984 kW/4,0461 kVA antes de arranques. La suma incluye ambos bancos y no se multiplica por dos otra vez al pasar a un solo camino. Su caudal efectivo, simultaneidad, protección y aire de reposición corresponden a ENV-19 \parencite{lote19-systemair-ex180}.

Sobre la envolvente intermedia de ECI-4, ya con margen propio del veinte por ciento, la extracción agrega 1,1904 kW/2,4276 kVA. Resultan \emph{50,3600 kW/54,4325 kVA parciales}, cociente de cotas $P/S=0,9252$, antes de unidades de tratamiento, bombas, frío, recalentamiento, controles y sus pérdidas. Los márgenes globales de placa son 9,6400 kW y 5,5675 kVA; no son capacidad térmica ni garantía de arranque.

\begin{table}[htbp]
\centering\small
\caption{Cotas parciales por fase tras extracción, camino superviviente}\label{tab:sd4-fases-eci20}
\begin{tabularx}{\linewidth}{X R{25mm} R{25mm} R{25mm}}
\hline
\EncabezadoTabla{Partida con margen 20\,\%} & \EncabezadoTabla{L1 kVA} & \EncabezadoTabla{L2 kVA} & \EncabezadoTabla{L3 kVA} \\ \hline
ECI-4, antes de extracción & 19,3965 & 15,1163 & 17,4919 \\ \hline
Cuatro motores activos & 0,8092 & 0,8092 & 0,8092 \\ \hline
Total parcial & 20,2057 & 15,9255 & 18,3011 \\ \hline
\end{tabularx}
\FuenteTabla{Cálculo de Synaptix desde corriente de ENV-19. Redondeos heredados de fases difieren de la suma agregada sin redondear en 0,000112 kVA.}
\end{table}

L1 excede los 20 kVA antes de añadir tratamiento. A la tensión de comprobación de 225,4 V, los 84,511 A anteriores más $1,20(4)(0,73)=3,504$ A resultan 88,015 A, superiores a los 86 A publicados de la UPS. Trasladar los seis juegos móviles de L1 a L2 quitaría sólo $1,20(6)(0,050/0,75)=0,480$ kVA: deja 85,8855 A, y el solape de un solo motor durante relevo agrega 0,876 A, hasta 86,7615 A, incluso antes de su arranque. Por ello no se adopta ese traslado como solución ni se reduce la cobertura de equipos para aparentar margen. Se necesita redimensionar o redistribuir el conjunto climático y su fuente con estados de arranque/relevo; conservar 60 kVA y el reparto actual no es una combinación admisible completa.

La demanda de supervivencia de cuatro motores es distinta del régimen normal A/B y del solape. Los cuatro motores activos suman 8.689,92 kWh/año si la ventilación es permanente durante las 8.760 h; cada hora adicional con otros cuatro añade 0,992 kWh. No se computa 8.689,92 por camino ni se presenta el margen de compra como consumo. Las reservas detenidas no se consideran sin consumo de control hasta individualizarlo. En batería también se conserva extracción; conexión previa a UPS sin otra fuente validada incumpliría esa frontera.

Para cribado de recarga con esa extracción y todavía sin tratamiento, $[0,94(0,99)(74,824595)-50,360015]/0,588=32,7751$ A agregados. Es techo aritmético bajo los supuestos anteriores; no es consigna Eaton disponible ni garantiza carga de 32,7751 A. Los límites de corriente de cadena y de fase se aplican además del límite agregado, y el tratamiento pendiente consume parte de ese margen.

\subsubsection{Generación: carga total, escalones y energía (GCI-20)}
\label{sec:sd4-escalones-gci20}

GCI-4 conserva la selección Avtron 4100/400 V/50 Hz, 60 kW y ALC para cada P200-6. La incompatibilidad documental de su auxiliar 120 V/60 Hz y la potencia del ventilador permanecen identificadas. Una especificación de bancada 3010 con motor de 50 Hz sería otro conjunto, con instalación interior y control SIGMA LT; no se trasladan sus 0,48 kW a la 4100 ni se atribuye ALC a otra serie. Se mantiene una sola selección vigente hasta disponer de una variante integral autorizada \parencites{lote17-avtron-4100}{lote17-avtron-alc}.

Para cada instante y tensión medida, el permiso de bancada comprueba
\[
S_{serv}(s)+S_{aux}(s)+S_{bancada}(s)\leq90\ \mathrm{kVA},\qquad
P_{serv}(s)+P_{aux}(s)+P_{bancada}(s)\leq90\ \mathrm{kW}.
\]
También comprueba corriente de cada fase, protección de 160 A y respuesta del alternador; $90$ kVA equivale a 129,9038 A sólo para carga trifásica equilibrada. La tolerancia superior de un escalón nominal de 5 kW a 400 V es 5,25 kW; con variación de tensión su cota resistiva se corrige por $(U/400)^2$. Antes de incorporar un escalón se incluye su cota superior completa y el arranque de su ventilador, no sólo su resolución nominal. Con servicio en la cota de 108 A/400 V, quedan 15,1754 kVA para bancada y auxiliares en conjunto. Un ejemplo de reserva externa propia de 3 kVA deja 12,1754 kVA: permite como máximo dos escalones de 5,25 kW si realmente están individualizados y sin agregar otro arranque; no habilita 60 kW de bancada. Esa reserva de 3 kVA es un presupuesto de comprobación, no el consumo de un ventilador aún sin ficha específica. El ALC tiene retardos y mide servicio, por lo que el permiso total y su retirada prioritaria son funciones independientes \parencite{lote17-avtron-alc}.

El retorno incorpora primero la extracción y seguridad sostenidas, después el servicio y tratamiento compatibles, y por último recuperación del banco y carga resistiva permitidas por el total. No se aplica una rampa ajena a esta UPS ni se presume que el ALC retire carga antes de cualquier escalón. La prueba registrará frecuencia, tensión mínima/máxima, duración fuera de banda, corriente por fase, DC de banco y comportamiento de receptores; las ventanas admisibles se fijan desde cada equipo seleccionado. La clase ISO 8528 G2 del grupo no constituye esa respuesta del conjunto \parencite{p7-fg-p200-ficha}.

La ficha publica 23,7 L/h a 50\,\% PRP, 32,6 L/h a 75\,\% y 41,3 L/h a 100\,\%, para 50 Hz en condiciones de referencia. No se extrapola 23,7 L/h hacia demanda inferior ni se interpreta el punto de 50\,\% como instrucción de mantenimiento. El objetivo propio de 57,6 kW de GCI-4 tampoco es un umbral FG Wilson. Se conserva la reserva individual de 1.500 L útiles y su límite de 52,0833 L/h para 24 h con veinte por ciento de margen; potencia/consumo corregidos, régimen de baja carga y transitorios deben corresponder al grupo, admisión y bancada seleccionados \parencite[pp.~1--3]{p7-fg-p200-ficha}.

\textbf{Frontera anual verificable.} Por estado excluyente $s$ de red, batería, grupo, retorno/recarga y prueba, se suman horas hasta 8.760 y se registra energía eléctrica importada, generación eléctrica local, exportación a puestos/campo, TI interior, auxiliares y variación de energía almacenada. La energía eléctrica del recinto es importación eléctrica más generación eléctrica, menos exportación y variación neta de almacenamiento dentro de esa frontera. Los litros o energía química de combustible no se añaden a ese numerador eléctrico. En un ciclo inicial/final equivalente, la variación de almacenamiento es cero; las pérdidas permanecen dentro. El PUE proyectado requiere el mismo período y frontera para numerador y TI interior. La energía útil de descarga ya alimenta las cargas; sumar nuevamente su recarga bruta duplica servicio. Los ensayos de bancada exteriores son carga externa, con sus pérdidas de suministro separadas. La adición de 8.689,92 kWh de ventilación seleccionada no basta para recalcular un PUE completo sin electricidad del tratamiento y curvas por estado: los valores antiguos 4,74/4,78 continúan como referencias condicionadas.
